% Source: source/普通高考/2011/2011浙江理.pdf, page 1, question 12.
% Type: program flowchart; pdfimages -list reports no embedded bitmap (vector line art).
% Grid evidence: original page rendered at 180 dpi; grid step=100, fine_step=50.
% Node-edge list: 开始→k=2→k=k+1→a=4^k→b=k^4→a>b；是→输出k→结束；
% 否→右侧回边→k=k+1入口主干。All borders/connectors are solid; labels are
% beside exits. Nodes are centered with local parindent=0pt; loop and main arrows
% remain distinct at the shared stem.
\begin{tikzpicture}[
  x=1cm,
  y=0.9cm,
  >=Stealth,
  line width=0.45pt,
  line cap=round,
  line join=round,
  every node/.style={font=\normalsize, inner sep=1pt, align=center, anchor=center,
    execute at begin node=\setlength{\parindent}{0pt}},
  flow/.style={draw, line width=0.45pt},
  terminal/.style={flow, rounded corners=3pt, minimum width=1.25cm, minimum height=0.70cm},
  process/.style={flow, minimum height=0.72cm, align=center},
  io/.style={flow, trapezium, trapezium left angle=74, trapezium right angle=106,
    minimum height=0.76cm, align=center},
  decision/.style={flow, diamond, aspect=2.9, minimum width=3.35cm,
    minimum height=1.0cm, inner sep=1pt, align=center}
]
  \node[terminal] (start) at (0,9.55) {开始};
  \node[process, minimum width=1.55cm] (init) at (0,8.15) {$k=2$};
  \node[process, minimum width=2.55cm] (increment) at (0,6.70) {$k=k+1$};
  \node[process, minimum width=2.05cm] (apow) at (0,5.30) {$a=4^k$};
  \node[process, minimum width=2.05cm] (bpow) at (0,3.90) {$b=k^4$};
  \node[decision] (test) at (0,2.45) {$a>b?$};
  \node[io, minimum width=2.20cm] (output) at (0,1.05) {输出 $k$};
  \node[terminal] (finish) at (0,-0.30) {结束};

  % Place the return segment below init with at least 2.5 mm clearance.
  \coordinate (loopJoin) at (0,7.40);
  \coordinate (loopRight) at (2.55,7.40);

  % Main path and terminating branch.
  \draw[->] (start.south) -- (init.north);
  \draw[->] (init.south) -- (increment.north);
  \draw[->] (increment.south) -- (apow.north);
  \draw[->] (apow.south) -- (bpow.north);
  \draw[->] (bpow.south) -- (test.north);
  \draw[->] (test.south) -- node[right=2pt] {是} (output.north);
  \draw[->] (output.south) -- (finish.north);

  % 否 branch: right and upward, then a single left-pointing arrow into the
  % shared stem before k=k+1; the main stem retains its own downward arrow.
  \coordinate (loopSide) at (2.55,2.45);
  \draw[->] (test.east) -- node[above=3pt] {否} (loopSide) -- (loopRight) -- (loopJoin);
\end{tikzpicture}
